\section{Memory requirements and performance models}
\label{sec:performance-background}
\label{sec:memory-requirements}

\subsection{Weight, KV and workspace capacity}

\tabref{tab:inference-data} compares the main data classes.  Their
lifetimes explain why process memory does not simply follow a request's
context length.  Persistent weights, reusable KV pools and temporary
workspaces can coexist even when only part of the pool is populated.

\begin{table}[htbp]
\centering
\small
\tabsetup
\begin{tabularx}{\textwidth}{P{2.1cm}YYY}
\headrow
\headcell{Property} & \headcell{Weights} & \headcell{KV state} & \headcell{Temporary data} \\[3pt]
\tabheadrule
\tabstripeRow
Origin & Learned parameters & Per-layer keys and values computed from inputs & Intermediate results and kernel scratch space \\[3pt]
\tabrowsep
\tabplainRow
Lifetime & Model residency, often many requests & Needed history; pool can outlive requests & Operation or phase; allocation can be reused \\[3pt]
\tabrowsep
\tabstripeRow
Modification & Normally read-only after preparation & New entries written; history read & Written and consumed during computation \\[3pt]
\tabrowsep
\tabplainRow
Sharing & Shared by requests using the same representation & Per sequence logically; compatible prefixes can share blocks & Usually tied to current execution; buffers can be reused \\[3pt]
\tabrowsep
\tabstripeRow
Size drivers & Parameter count, types, format, packing, copies & Layers, KV heads, head size, retained positions, type, pool policy & Batch, prompt chunk, operator, implementation \\[3pt]
\tabrowsep
\tabplainRow
Access & Dense transforms reuse matrix tiles; embeddings select rows & Append state and read permitted history & Short-lived reads and writes; fusion can reduce materialization \\[3pt]
\tabbottomrule
\end{tabularx}
\caption{Origins, storage roles and access behavior of inference data.
A logical lifetime does not imply immediate deallocation or page eviction.}
\label{tab:inference-data}
\end{table}
\FloatBarrier

For \(N\) distinct stored parameters, uniform payload size \(s_W\)
bytes per parameter, and representation overhead \(M_{\mathrm{format}}\),
a weight representation occupies approximately
\begin{equation}
\label{eq:weight-payload}
 M_W \approx Ns_W + M_{\mathrm{format}}.
\end{equation}
\eqcaption{Weight representation estimate.  Format overhead includes
quantization metadata, alignment and padding; additional runtime copies
are accounted for separately.}

With mixed tensor types, the payload term becomes
\(\sum_j N_j s_j\), counting shared storage once.  Packing overhead must
be included in the representation it affects.  A training checkpoint
may additionally contain optimizer state that inference does not need,
so checkpoint bytes need not equal the inference weight payload.

As an illustrative example, exactly \(8\times10^9\) parameters stored
in BF16 require \(16\times10^9\) payload bytes: 16\,GB or approximately
14.90\,GiB.  Four-bit storage gives an ideal payload of 4\,GB
(3.73\,GiB).  If that illustrative format also stores one two-byte scale
per 32 weights, the scales add 0.5\,GB, giving 4.5\,GB (4.19\,GiB)
before other metadata and padding.  This is an example format, not a
size claim for every four-bit model.  A retained second copy of the
BF16 representation would add another 16\,GB regardless of the file's
size.  Here GB denotes \(10^9\) bytes and GiB denotes \(2^{30}\) bytes;
KiB similarly denotes \(2^{10}\) bytes.

For \(L\) layers with uniform full-history attention, \(n_{\mathrm{KV}}\)
KV heads per layer, head dimension \(d_h\), and \(s_{\mathrm{KV}}\)
bytes per stored scalar, the live unshared KV payload is
\begin{equation}
\label{eq:kv-live-payload}
 M_{\mathrm{KV,live}}=
 2L n_{\mathrm{KV}}d_h s_{\mathrm{KV}}\sum_i S_i.
\end{equation}
\eqcaption{Live KV payload for independent, unshared sequences.  The
factor two counts keys and values; \(S_i\) is the number of evaluated
positions whose state is retained for sequence \(i\).}

The dimensions are layers times heads times scalars per head times
bytes per scalar times positions, hence bytes.  \(S_i\) counts cached
positions, not automatically every emitted token, since the final
output may not have been evaluated.  Nonuniform layers require summing
per-layer payloads; bounded histories, shared prefixes, quantization
metadata and padding require corresponding adjustments.  Occupied block
capacity rounds each independently allocated length to block boundaries.
Reserved pool capacity can be larger still.

For an illustrative 32-layer model with eight KV heads, \(d_h=128\)
and two-byte entries, each retained position occupies
\(2\times32\times8\times128\times2=131{,}072\) bytes, or 128\,KiB.
An 8{,}192-position history therefore holds 1\,GiB of live KV payload.
Thirty-two independent sequences of that length hold 32\,GiB.
Replacing eight KV heads with 32, while keeping the other assumptions,
multiplies these sizes by four; using one KV head divides them by eight.
The examples isolate architecture and workload effects, without assuming
a particular named model or allocator.

A capacity budget must include allocations that coexist:
\begin{equation}
\label{eq:inference-allocation-budget}
 M_{\mathrm{allocated}} \approx M_W+M_{\mathrm{copies}}+
 M_{\mathrm{KV,allocated}}+M_{\mathrm{workspace}}+M_{\mathrm{other}}.
\end{equation}
\eqcaption{A simultaneous allocation budget.  KV allocation includes
its allocated pool or buffers, so live KV payload must not be added a
second time.  Other storage includes runtime metadata and allocator overhead.}

Combining the illustrative BF16 weights with the 32 unshared histories
gives about 46.90\,GiB of weight and live-KV payload before workspace and
allocation overhead.  A reserved pool larger than the live state raises
the allocation budget.  Temporary conversion buffers can instead make
loading the peak-memory interval.  The sum describes allocated storage,
not necessarily resident pages or controller traffic; neither logical
payload formulas nor process RSS alone identify all these components.

\subsection{Reuse and access traffic}

For the dense operation in \eqnref{eq:weights-activations}, there are
approximately \(2m d_{\mathrm{in}}d_{\mathrm{out}}\) floating-point
operations.  If its weight matrix is fetched once across the memory
boundary, weight traffic is \(s_W d_{\mathrm{in}}d_{\mathrm{out}}\)
bytes.  Neglecting all other traffic then gives the illustrative ratio
\(2m/s_W\) FLOP/byte.  It is a weight-only estimate: input and output
traffic, packing, cache reuse between evaluations and repeated tile
loads determine the actual arithmetic intensity.  Quantized kernels
also require a suitable operation-count convention.

Prefill and batches with several evaluated positions can reuse a weight
tile across activation rows.  Low-batch decode offers less of this
reuse, while longer histories increase attention's KV reads.  Distinct
sequences generally access distinct histories, although grouped heads,
shared prefixes and kernel tiling can allow further reuse.  Temporary
buffers can contribute significant traffic when intermediate results
are materialized.  Neither weights nor KV state should therefore be
declared universally dominant: the relevant balance depends on the
model, representation, phase, workload and memory boundary.

An estimated logical read of \(D\) bytes is not a measurement of
main-memory traffic.  Cache hits, writeback, prefetching and transfers
between locality domains alter the observed byte counts.  If some data
resides in HBM and some in DDR, their separate traffic and available
paths must be considered.  A single bandwidth ceiling cannot simply
be applied to the sum of both tiers' capacities.  Simultaneous transfers
may overlap, but they can also share compute resources or interconnects.
Migration adds its own reads, writes, synchronization and possible
capacity peak; a benefit must cover these costs over the interval in
which the new placement is used.

\subsection{Compute throughput and practical ceilings}

\textbf{Compute throughput} counts operations per unit time.  Floating-point
rates use FLOP/s; a multiplication and addition count as two operations,
including when fused into one instruction.  Peak throughput depends on
active cores, clock frequency, supported instructions and precision.
It is distinct from inference throughput in tokens/s, since one token
requires many operations and accesses.

Peak specifications describe ideal conditions.  Microbenchmarks can
establish practical compute and bandwidth ceilings for a chosen
configuration, while application measurements show how much of that
capability is used~\cite{nerscroofline2026}.  The capacity, latency and
bandwidth distinctions in \secref{sec:cpu-memory-systems} remain important:
low bandwidth utilization alone does not establish a compute bottleneck.

\subsection{Bandwidth requires concurrency}
\label{sec:littles-law}

High bandwidth requires enough independent transfers to keep memory
busy.  Little's law relates sustained bandwidth \(B\), mean access
latency \(\tau\), and mean bytes in flight \(C\):
\begin{equation}
\label{eq:little}
 C=B\tau.
\end{equation}
\eqcaption{Mean memory concurrency required in steady state.  Bytes/s
multiplied by seconds gives bytes in flight.}

Dividing \(C\) by cache-line size gives the corresponding outstanding
line count.  Finite per-core tracking resources limit how many misses
can overlap.  More cores, independent software accesses and prefetching
can supply additional requests, subject to those limits.  Use loaded
latency: queueing near saturation can raise it above idle latency, so an
idle measurement gives only an optimistic concurrency requirement.
Regular streams can expose this parallelism; dependent load chains may
not.  \secref{sec:platform-concurrency} applies the relationship to the
experimental platform.  This memory-request concurrency differs from
both active inference requests and execution-batch size.

\subsection{Reading a Roofline plot}
\label{sec:roofline}

The Roofline model relates a kernel's operation count to its data
movement~\cite{williams2009roofline}.  Let \(F\) be the number of
floating-point operations and \(D\) the bytes read and written across
the memory boundary being analysed.  Their ratio is the
\textbf{arithmetic intensity}:

\begin{equation}
\label{eq:arithmetic-intensity}
  I = \frac{F}{D}
  \qquad \text{[FLOP/byte]}.
\end{equation}
\eqcaption{Arithmetic intensity: the amount of computation performed
per byte transferred.}

For a main-memory Roofline, \(D\) counts traffic between main memory and
the cache hierarchy, rather than every access issued by the processor.
Reusing data in cache can reduce this traffic and increase \(I\).
If the computation takes time \(t\), its measured compute throughput is
\(P=F/t\), and its achieved bandwidth at that boundary is \(D/t\).

With a compute ceiling \(P_{\max}\) and bandwidth ceiling
\(B_{\max}\), the model gives

\begin{equation}
\label{eq:roofline}
  P \leq \min\!\left(P_{\max},\, B_{\max} I\right).
\end{equation}
\eqcaption{The Roofline bound: attainable throughput is limited by
both the processor's arithmetic rate and the supply of data.}

\figref{fig:roofline} plots this bound on logarithmic axes.  The
sloping part is the bandwidth limit \(B_{\max}I\); the horizontal
part is the compute limit \(P_{\max}\).  They meet at the
\textbf{ridge point}, whose intensity is
\(I_{\mathrm{ridge}}=P_{\max}/B_{\max}\).

\begin{figure}[htbp]
\centering
\begingroup
\definecolor{roofink}{HTML}{263445}
\definecolor{roofblue}{HTML}{365D85}
\colorlet{roofaccent}{sapred}
\definecolor{roofborder}{HTML}{CBD5DF}
\resizebox{\textwidth}{!}{%
\begin{tikzpicture}[
  font=\small\sffamily, text=roofink,
  axis/.style={-{Stealth[length=5pt,width=4pt]},
               draw=roofink!80, line width=0.8pt},
  note/.style={font=\footnotesize\sffamily, align=center},
  card/.style={draw=roofborder, fill=roofink!3, rounded corners=4pt,
               line width=0.65pt, inner xsep=7pt, inner ysep=7pt,
               text width=4.3cm, minimum height=2.5cm, align=left,
               font=\footnotesize\sffamily}
]
  % x = 3 log10(I / 0.1); y = 1.5 log10(P / 10).
  % Illustrative ceilings: 100 GB/s and 1000 GFLOP/s; ridge at 10 FLOP/byte.
  \fill[roofblue!7] (0,0) -- (6,3) -- (6,0) -- cycle;
  \fill[roofaccent!7] (6,0) rectangle (9,3);
  \foreach \x in {3,6,9}
    \draw[roofborder, line width=0.4pt] (\x,0) -- (\x,3.35);
  \foreach \y in {1.5,3}
    \draw[roofborder, line width=0.4pt] (0,\y) -- (9,\y);

  \draw[axis] (0,0) -- (9.45,0);
  \draw[axis] (0,0) -- (0,3.8);
  \foreach \x/\value in {0/0.1,3/1,6/10,9/100}
    \draw[roofink!70] (\x,0) -- (\x,-0.08)
      node[below, font=\footnotesize\sffamily] {\value};
  \foreach \y/\value in {0/10,1.5/100,3/{1{,}000}}
    \draw[roofink!70] (0,\y) -- (-0.08,\y)
      node[left, font=\footnotesize\sffamily] {$\value$};
  \node[note, rotate=90] at (-1.15,1.75)
    {Compute throughput $P$ [GFLOP/s]};
  \node[note] at (4.5,-0.65)
    {Arithmetic intensity $I$ [FLOP/byte]};

  \draw[roofaccent!55, densely dashed, line width=0.7pt]
    (0,3) -- (6,3);
  \draw[roofblue, line width=1.9pt] (0,0) -- (6,3)
    node[pos=0.58, sloped, above=4pt, note, text=roofblue]
      {Bandwidth limit: $B_{\max} I$};
  \draw[roofaccent, line width=1.9pt] (6,3) -- (9,3)
    node[midway, above=4pt, note, text=roofaccent]
      {Compute limit: $P_{\max}$};

  \draw[roofink!60, densely dashed] (6,0) -- (6,3);
  \fill[roofink] (6,3) circle (2.3pt);
  \node[note, anchor=south] (ridge) at (5.3,3.6)
    {Ridge point: $I_{\mathrm{ridge}}=P_{\max}/B_{\max}$};
  \draw[roofink!65, -{Stealth[length=4pt,width=3pt]}]
    (ridge.south) -- (6,3.08);

  \node[note, text=roofblue, font=\small\sffamily\bfseries]
    at (2.65,4.35) {Memory-bound region};
  \node[note, text=roofaccent, font=\small\sffamily\bfseries]
    at (7.6,4.35) {Compute-bound region};

  % Points are illustrative measurements below the roof.
  \draw[roofink!60, densely dashed] (3,1.05) -- (3,1.5);
  \fill[roofink] (3,1.05) circle (2.6pt);
  \node[note, anchor=north west] at (3.1,1.0) {Kernel A};
  \draw[roofink!60, densely dashed] (7.8,2.6) -- (7.8,3);
  \fill[roofink] (7.8,2.6) circle (2.6pt);
  \node[note, anchor=north] at (7.8,2.45) {Kernel B};

  \node[card, draw=roofblue!45, fill=roofblue!4] at (2.2,-2.3)
    {\textbf{Left of the ridge}\\
     Limited by data movement.\\
     More bandwidth lifts the roof.\\
     Reuse moves a kernel right.};
  \node[card, draw=roofaccent!45, fill=roofaccent!4] at (7.2,-2.3)
    {\textbf{Right of the ridge}\\
     Limited by arithmetic.\\
     Better use of compute units\\
     brings a kernel closer to the roof.};
\end{tikzpicture}%
}
\endgroup
\caption{A Roofline plot with illustrative ceilings of 100\,GB/s and
1{,}000\,GFLOP/s, giving a ridge at 10\,FLOP/byte.  Both axes are
logarithmic.  The sloping segment is the bandwidth limit and the
horizontal segment is the compute limit.  The roof is an upper
bound; the dots represent example kernel measurements below it, and the
dashed vertical segments show their gap to the bound.  This is a
conceptual example, not a measurement of the thesis platform.}
\label{fig:roofline}
\end{figure}

Left of the ridge, the model identifies a \textbf{memory-bound}, more
specifically bandwidth-bound, region.  Increasing bandwidth raises the
bound, while reusing each transferred byte more often moves a kernel to
the right.  Right of the ridge, the bound is \textbf{compute-bound};
additional bandwidth alone does not raise it.  For example, with the
figure's ceilings, intensity \(1\) limits throughput to 100\,GFLOP/s,
whereas intensity \(20\) permits at most 1{,}000\,GFLOP/s.

A point below the roof does not establish the cause of the gap.
Latency, limited parallelism, inefficient instructions and runtime overhead
can all reduce performance.  Compare kernels with ceilings for the same
cores, precision and memory boundary~\cite{nerscroofline2026}; integer
kernels require matching integer-operation counts and ceilings, not FLOP/s.

Prefill's weight reuse and low-batch decode's smaller reuse motivate
separate phase analysis.  Their actual intensity depends on batching,
history length, representation and kernels~\cite{na2024llmcpu}.
Roofline describes numerical work; request timing also includes scheduling
and other runtime overhead.
\FloatBarrier
